\section{Future work}
\label{sec:future}

\paragraph{Run the registered experiment.}
The immediate next step is the one this paper is missing: execute the frozen
protocol at $256$~px and report the verdict, whichever way it falls. The
estimated cost is about one GPU-hour. Everything needed is written and tested.

\paragraph{Compression and detection latency (drafted, deliberately not run).}
The natural follow-up becomes available only once a monitoring harness exists.
\citet{fu2026quantization} report that post-training quantization of a world
model produces \emph{asymmetric} sensitivity across components, and failure
modes that further optimization does not repair. A mean-performance number hides
that; a per-anomaly-family detection-latency measurement should not.

We have drafted a second study asking whether compression changes \emph{which}
anomalies a model still catches and \emph{how fast}, rather than degrading it
uniformly. Its shape:

\begin{itemize}
\item Six precisions (fp32, bf16, fp16, int8 weight-only, W8A8, 4-bit
weight-only) crossed with three placements (encoder only, predictor only,
both). Placement is the factor that tests the asymmetry claim directly: under
encoder-only compression the latents the probes consume are degraded while the
prediction head is intact, and vice versa.
\item One primary contrast per anomaly family, Holm-corrected across the four,
with every other cell exploratory and labelled as such.
\item \textbf{Equivalence testing}, not null-hypothesis testing. The interesting
claim is usually that compression does \emph{not} hurt, which is an equivalence
claim and cannot be supported by failing to reject a difference. The draft
specifies TOST on the latency difference --- and leaves the equivalence margins
blank, as an explicit TODO for the author to set and justify before freezing,
because choosing a margin after seeing the effect is the easiest way to
manufacture an equivalence result.
\item A fresh seed for the confirmatory clips, so that no clip used to develop
the harness appears in the confirmatory set, and per-cell threshold calibration,
so that a shift in score scale is not confounded with a change in sensitivity.
\item A second-GPU arm. Quantization kernels, tensor-core paths and accumulation
order differ across hardware, so a compression effect that does not reproduce on
two GPUs is a kernel artefact rather than a property of the model.
\item Two exploratory analyses: \emph{blind spots} --- whether misses
concentrate in one anomaly family rather than thinning uniformly --- and
\emph{motion energy}, testing whether compression costs concentrate in
low-motion clips where there is least signal to work with. Both target the same
idea, that aggregate accuracy can hold while the distribution of failures
shifts.
\end{itemize}

The draft is marked \textsc{draft, not frozen} and no data may be collected
under it until its margins are set and it is frozen in its own commit. It does
not amend the v1 pre-registration, which remains in force untouched. We record
it here so that its existence and date are on the record and it cannot later be
presented as having been registered earlier than it was.

\paragraph{Other directions.}
A physics-simulator domain with contact and articulation would test whether
anything here survives outside synthetic discs. A gradual-onset benchmark would
let the anticipation question be asked properly, rather than the latency
question we can currently answer. And the per-clip standardisation issue raised
in Section~\ref{sec:limitations} deserves a clean treatment: a probe family
defined to be invariant to per-clip offsets from the start, rather than one
patched after the fact.
